% Generated draft: requires mathematical and source review.
\documentclass[11pt]{article}
\usepackage[margin=25mm]{geometry}
\usepackage{fontspec}
\setmainfont{Latin Modern Roman}
\usepackage{amsmath,amssymb,mathtools,unicode-math}
\setmathfont{Latin Modern Math}
\usepackage{enumitem,needspace,xurl,hyperref}
\setlength{\parindent}{0pt}
\setlength{\parskip}{5pt}
\setlength{\emergencystretch}{3em}
\newcommand{\N}{\mathbb{N}}
\newcommand{\Z}{\mathbb{Z}}
\newcommand{\Q}{\mathbb{Q}}
\newcommand{\R}{\mathbb{R}}
\newcommand{\C}{\mathbb{C}}
\newcommand{\F}{\mathbb{F}}
\newcommand{\A}{\mathbb{A}}
\newcommand{\PP}{\mathbb{P}}
\newcommand{\E}{\mathbb{E}}
\newcommand{\eps}{\varepsilon}
\newcommand{\dd}{\,\mathrm{d}}
\begin{document}
\section*{UnsolvedMath ID: 10300051}
\textbf{Source problem code:} catalogue entry
\subsection{Definitions and mathematical statement}
Let \(\Gamma=\pi_1(M)\), with a word metric from a finite generating set when available. Let \(\rho:\Gamma\to\R\) be a quasimorphism: there exists \(C<\infty\) with
\[
 |\rho(\alpha)+\rho(\beta)-\rho(\alpha\beta)|<C
 \qquad(\alpha,\beta\in\Gamma).
\]
No homogeneity condition \(\rho(g^n)=n\rho(g)\) is stated. For the interval \(I\) used in the source, set \(L_\rho=\rho^{-1}(I)\), with the metric inherited from \(\Gamma\). The interval's location, width and relation to \(C\) are absent from the excerpt.

For a fixed scale \(R\), coarse connected components can be interpreted as equivalence classes under chains of points in \(L_\rho\) whose consecutive word distances are at most \(R\). A usual coarse simple-connectivity formulation uses Rips complexes: \(P_r(A)\) has vertex set \(A\) and a simplex for each finite set of diameter at most \(r\); after coarse connectedness, one asks that for every \(r\) there be \(R\geq r\) such that loops in \(P_r(A)\) become null-homotopic in \(P_R(A)\). The source does not fix its component scale, Rips convention or uniformity requirements.

Question 12.1 asks:
\begin{enumerate}
\item Are the coarse connected components of \(L_\rho\) coarsely simply connected?
\item If yes, is there a meaningful sense in which \(L_\rho\) and its translates \(gL_\rho=\{gx:x\in L_\rho\}\), \(g\in\Gamma\), form coarse minimal planes in \(\Gamma\)?
\end{enumerate}
The second part asks for a notion as well as its validity; ``coarse minimal plane'' is not defined and must not be replaced silently by a quasi-isometrically embedded plane or a particular area-minimizing condition. Since \(\rho\) is not necessarily a homomorphism, translates of a level band are not automatically exact level bands at shifted values. The intended 3-manifold hypotheses on \(M\), finite generation, and interval restrictions need source context.
\subsection{Short English statement}
Are coarse components of a bounded-defect cochain's level band coarsely simply connected? If so, can those bands and their group translates be interpreted as coarse minimal planes? The interval, manifold class, and coarse conventions are missing.
\subsection{Sources}
\begin{itemize}
\item Local UnsolvedMath registry, permanent ID 10300051, source code catalogue entry; no individual proposer or proposal year supplied; update 2026-08-06.
\item Calegari, \emph{Problems in foliations and laminations of 3-manifolds} (2002), Question 12.1: \url{https://arxiv.org/abs/math/0209081}. Registry access date 2026-07-29; source-TeX extraction. The background reports 67 numbered questions. Current primary-literature status and redistribution terms require review.
\item The supplied summary mentions bounded cohomology and quasimorphism level sets without exact further references or a verified answer.
\end{itemize}
\end{document}
